\documentclass[a4paper,11pt]{ltjsarticle}
\usepackage[margin=20mm]{geometry}
\usepackage{booktabs,tabularx,array,enumitem,xcolor}
\usepackage{tikz}
\usepackage[most]{tcolorbox}
\usetikzlibrary{positioning,arrows.meta,calc}
\usepackage[hidelinks]{hyperref}
\setlist{nosep,leftmargin=1.5em}
\renewcommand{\arraystretch}{1.25}
\newcolumntype{L}{>{\raggedright\arraybackslash}X}
\definecolor{acc}{HTML}{1F4E79}
\definecolor{soft}{HTML}{EAF1F8}
\definecolor{warn}{HTML}{FFF4E5}
\tcbset{colback=warn,colframe=acc,boxrule=0.6pt,arc=2pt}
\newcommand{\note}[1]{{\small\color{black!65}#1}}

\title{DAI 論文の現状\\[2pt]\large 論理の流れ，先行研究との違い，実験の不備，比較の不足，計算時間，削減案}
\author{MER-DA（作成：Claude．研究者への報告）}
\date{2026 年 11 月 1 日}

\begin{document}
\maketitle

\section*{要点}

\begin{enumerate}
\item 論文は 6 つの筋で組む．ソースを混ぜたときの得と損を測り（筋 1），得の理由を分類器の 2 種類の誤りで説明し（筋 2〜4），その説明から作った手法（筋 5）が標準の手法を上回る理由を部品ごとに示す（筋 6）．
\item 先行研究との違いは，部品ではなく組み合わせにある．中心合わせ，ソースの分類器を出発点にした再学習，別のモダリティの分類器からの蒸留，ソースを仮のターゲットにした設定の選び方は，どれも先例がある．57 本の論文を読んだ範囲で，これらを凍結した特徴の上の線形の分類器で同時に使い，K を 0〜50 で振って測った研究は無かった．
\item TAC 版の模擬査読の指摘を筋ごとに数えた 58 件のうち，27 件は新しい計算が要る状態だった．今朝の E75 で保存済みの結果から出せるものを計算し，筋 1〜3 の言い方を 4 か所弱めた．
\item 今ある比較（E61）には，標準のドメイン適応の手法（DANN，MME，CORAL など）と，提案と同じ参照分類器を使う蒸留が無かった．E77・E78 で足す．評価に使うターゲットはすべて手法の設計の段で見ているので，「見ていないコーパスで確かめていない」という不備は計算では直らない．
\item 今の設定のまま計算すると，判定の解析（E79）を回せるのは 11 月 7〜16 日ごろになる（E78 の所要をまだ正常な GPU の状態で測れていないので幅が広い）．7 節の削減案を入れると，E77 が約 1 日早く終わる．
\end{enumerate}

\section{この資料で使う言葉}

\begin{tabularx}{\linewidth}{@{}p{0.25\linewidth}L@{}}
\toprule
言葉 & 意味 \\
\midrule
ターゲット & 感情を認識したい新しいコーパス．ラベルは感情ごとに K 本しか無い \\
ソース & ラベルが全部付いた別のコーパス \\
UAR & 感情ごとの正解率の平均．4 感情なので，でたらめに答えると 25\% \\
合併 & ソースの全件とターゲットの K 本を一緒に学ぶこと．ターゲットのみ，ソースのみと比べる \\
利得 & 合併の UAR からターゲットのみの UAR を引いた値．単位はポイント（点） \\
o（判断基準のずれ） & 分類器の切片（各感情の点数に足す定数）がずれていることによる誤り．切片を直せば消える \\
d（識別力） & 分類器の重みの向きで決まる，感情どうしの見分けやすさ．向きを変えないと直らない \\
向きの損 & d が最良の向きより劣るために失う正解率 \\
中心合わせ & ソースとターゲットの特徴から，それぞれの平均を引く操作．ラベルを使わない \\
仮説転移 & ソースで学んだ分類器を出発点にして，ターゲットの K 本で学び直すこと \\
出発点（E39） & 全ソースを混ぜ，5 つの設定値を仮のターゲットで選んで作った分類器 \\
候補 & ラベルの無いターゲットの発話（テストの話者を除く） \\
弱いラベル & 候補に付ける，別の分類器が出した確率の答え \\
参照分類器 & 弱いラベルを出す分類器 \\
ours，oursx & 提案手法．ours は別のモダリティの分類器の答えを弱いラベルにする．oursx は参照分類器・重みの付け方・弱いラベルの総重みを候補にして，K ごとに選ぶ \\
仮のターゲット & ソースのコーパスの 1 つをターゲットの代わりにして設定値を選ぶ手順．テストのラベルを使わずに設定を決めるため \\
世界 & 使うモダリティの組．AVT は音声・顔・言葉，AV は音声・顔 \\
生徒 & 学び直す側の分類器．融合の生徒は 3 つ（AVT）か 2 つ（AV）の特徴をつないで入力にする \\
\bottomrule
\end{tabularx}

\section{論文の論理の流れ}

図 1 に 6 つの筋のつながりを示す．各筋の答えが次の筋の問いを生む順に並べた．

\begin{figure}[h]
\centering
\begin{tikzpicture}[
  box/.style={draw=acc,fill=soft,rounded corners=2pt,text width=0.78\linewidth,inner sep=5pt,align=left,font=\small},
  lab/.style={font=\small\color{acc},fill=white,inner sep=1pt},
  arr/.style={-{Stealth[length=2.2mm]},thick,acc},
  node distance=7mm]
\node[box] (s1) {\textbf{筋 1　利得の地図}\quad ソースを混ぜると UAR は上がるか．\\
→ K が小さいと上がり（K=1 で +7.3〜+12.9 点），K=25 以上で下がる．どの組で上がるかは，コーパス間の距離では説明できない．};
\node[box,below=of s1] (s2) {\textbf{筋 2　判断基準のずれ o と識別力 d}\quad 少ないラベルは分類器の何を直すか．\\
→ 主に切片（o）を直す．向き（d）はほとんど直らない．o はラベル無しの中心合わせでも大部分が消える．};
\node[box,below=of s2] (s3) {\textbf{筋 3　ソースの価値}\quad 向き d はどこから補えるか．\\
→ ソースの分類器を出発点にすると向きを受け取れる．1 ソースで感情ごとに約 10 本のラベルにあたる．};
\node[box,below=of s3] (s4) {\textbf{筋 4　ターゲットの価値}\quad 出発点に残る誤りは何か．ターゲットのデータで直るか．\\
→ 残るのは向き d の分．K 本のラベルでもラベルの無い発話でも向きは決まらない．};
\node[box,below=of s4] (s5) {\textbf{筋 5　蒸留}\quad 向きを動かす情報はどこにあるか．\\
→ 別のモダリティの分類器の答え．候補の全部に弱いラベルとして付ける（ours，oursx）．};
\node[box,below=of s5] (s6) {\textbf{筋 6　なぜ良いか}\quad 得はどの部品から来るか．\\
→ 部品を 1 つずつ外した実験と，全手法の向きの損の測定で分ける（E77〜E79．計算中）．};
\draw[arr] (s1) -- node[lab,right=2pt]{なぜ上がるのか} (s2);
\draw[arr] (s2) -- node[lab,right=2pt]{d はどう補うか} (s3);
\draw[arr] (s3) -- node[lab,right=2pt]{何が残るか} (s4);
\draw[arr] (s4) -- node[lab,right=2pt]{ラベル以外の情報は} (s5);
\draw[arr] (s5) -- node[lab,right=2pt]{なぜ上回るか} (s6);
\end{tikzpicture}
\caption{論文の 6 つの筋．各箱の 1 行目が問い，2 行目が今の答えである．}
\end{figure}

\subsection*{各筋の根拠と，今朝の E75 で変わった点}

表 1 に，各筋の主な数値と，E75 の再解析で言い方を変えた点をまとめた．数値は説明資料（\texttt{DAI\_paper/docs/sources/2026-10-31/explanation/p1〜p5}）と E75 の結果（\texttt{results/e75/}）から取った．

\begin{table}[h]
\centering\small
\caption{筋ごとの根拠の数値と E75 での変更}
\begin{tabularx}{\linewidth}{@{}p{0.09\linewidth}LL@{}}
\toprule
筋 & 根拠の数値（今の版） & E75 で変えた言い方 \\
\midrule
1 & 8 種類の特徴で K=1 の利得 +7.3〜+12.9 点．距離と利得の順位相関 $-$0.36〜+0.19 &
反復の間の依存を考えた区間にすると，K=1 で 0 より上の組は音声 A で 40 組中 31 組（素朴な区間では 38 組）．「ほぼすべての組で上がる」とは書かない．組の得の分散のうちコーパスの性質で説明できる割合は，調整すると A の K=5 で 0.36 で，「51〜73\%」とは書かない \\
2 & K=5 で，切片だけを学び直す手が合併の上がり幅の 0.91 倍（A），0.90 倍（A(HuBERT)），0.78 倍（V）を得る &
「K=5 の得の 5〜6 割は o の直し」という割合は根拠から外した．割合を出す模型（感情の対ごとに 2 つの山を仮定する式）の誤差が 0.52〜0.85 点で，測りたい位置の損と同じ大きさだった．代わりに左の UAR の単位の比で書く \\
2 & 式 (2) で，仮説転移の位置の損を予言する &
式 (2) は「模型で評価した位置の損」を 2 倍以内で当てるが，実際の正解率での損は予言の 3.9 倍だった．文を前者に限る \\
3 & 1 ソースの価値は感情ごとに約 10 本．出発点 E39 は K=5 で A 56.0（ターゲットのみ 42.0，合併 46.5） &
q は「有効次元」ではなく 1 つの数による要約と書く（A のターゲットごとの値の範囲 4.6〜90.7）．合併が o を直す割合 λ の式は，K ごとに組の並びを当てる（順位相関 A で 0.83〜0.87） \\
4 & 出発点と，候補の全ラベルで学べる最良との差は A で 11 点．K 本の選び方を変えても直らなかった（E33・E37・E40） & 変更なし．標準の DA の手法でも向きの損が減らないことを E77〜E79 で確かめる \\
5 & AVT の融合の生徒で ours は K=0 で +3.01，oursx は K=25 で +2.42．最も強い相手は K=0 で +1.44（massbal） & 変更なし．比較相手を足して E77〜E79 で測り直す \\
\bottomrule
\end{tabularx}
\end{table}

筋 2 の主張の書き直しと，夜の間に委任で決めたことは，研究者の承認を待っている（朝の報告の 2・4 節）．

\section{先行研究との違い}

\subsection{最も近い研究}

最も近いのは Feng \& Chaspari（IEEE TAC）である．演技の音声（IEMOCAP 台本，RAVDESS，CREMA-D，eNTERFACE）をソースに，IEMOCAP 即興をターゲットにして，話者・感情ごとに 1〜10 本のラベルで認識した．本研究とは次の点が違う．

\begin{tabularx}{\linewidth}{@{}p{0.22\linewidth}LL@{}}
\toprule
項目 & Feng \& Chaspari & 本研究 \\
\midrule
モダリティ & 音声だけ & 音声，顔，言葉とその組（8 種類） \\
特徴 & 手作りの 64 次元 & 学習しないエンコーダ（WavLM，ResNet，RoBERTa） \\
感情 & 3 つ & 4 つ \\
ターゲット & 1 つ & 7 群（組は最大 40） \\
ソース & 1 つずつ別々に & 1 つずつと，全部を混ぜる場合 \\
ずれの測定 & しない & 判断基準のずれ o と識別力 d に分けて，K の関数として測る \\
\bottomrule
\end{tabularx}

同論文は「複数のソースの組み合わせ」と「ソースの性質と少数ラベルでの有効性の関係」を今後の課題に挙げている．本研究はこの 2 つに答える．

\subsection{部品ごとの先例}

提案に使う部品は，1 つずつ見るとどれも先例がある（表 2）．

\begin{table}[h]
\centering\small
\caption{部品ごとの先例（\texttt{docs/related-work.md} 1.12〜1.15 節）}
\begin{tabularx}{\linewidth}{@{}p{0.33\linewidth}Lp{0.13\linewidth}@{}}
\toprule
部品 & 先例 & 判定 \\
\midrule
ソースの分類器を出発点にした再学習（L2 の罰則） & Kuzborskij \& Orabona 2013，Tommasi ほか 2010/2014，L2-SP 2018 & あり \\
固定した特徴の上の線形の分類器で，ソースの重みに罰則をかける半教師ありの DA & Donahue ほか 2013（ソースの超平面とのなす角に罰則．本研究は二乗距離） & あり \\
合併が切片を直す割合（λ） & De Silva 2023（1 次元の判別），Koh 2019 の命題 4 & 近い先例あり \\
別のモダリティの分類器から音声の生徒への蒸留（感情） & Albanie 2018（顔→音声），Srinivasan 2022，LiSER 2025 & あり \\
融合の出力を単モダリティの頭にまねさせる & xMUDA の Fusion，BFtD，FtD++ & あり \\
単モダリティの分類器の確率の積 & Kittler 1998 の積の規則ほか（最終の予測器として） & あり \\
弱いラベルの感情ごとの合計をそろえる & ReMixMatch の分布の整合，SeLa，SHOT の多様性の項，Saerens 2002 & 部分的 \\
ソースの 1 つを外して設定を選ぶ & Gulrajani \& Lopez-Paz 2021（leave-one-domain-out） & あり \\
\bottomrule
\end{tabularx}
\end{table}

\subsection{新しいと言える点}

57 本の本文を読んだ範囲（\texttt{docs/related-work.md} 1.14 節）で，次の組み合わせを同時に満たす研究は無かった．序論で「初めて」と書いてよいのは次の 4 点に限る．

\begin{enumerate}
\item 顔を含む 3 つのモダリティで，7 群の多数の組にわたり K を 0〜50 で振って利得を測った．
\item 利得を o と d に分けて K の関数として測り，中心合わせが置き換える分と同じ分け方で比べた．
\item ソースの価値をラベルの本数に換算した（1 ソースで感情ごとに約 10 本）．
\item ソースだけで学んだ別のモダリティの融合の分類器を，ターゲットで学び直さずに弱いラベルの元にする構成を，同じ条件で調整した標準の手法と K の全域で比べた．
\end{enumerate}

\subsection{書いてはいけないこと}

\begin{itemize}
\item 「融合の分類器から単モダリティへの蒸留は新しい」．xMUDA の Fusion と BFtD がある．残る違いは，参照分類器がソースで固定した線形の融合の分類器であること，運用時に単モダリティだけを使うこと，参照に確率の積を 2 段で使うことである．
\item 「新しい設定の選び方」．ソースの 1 つを外して選ぶ手順は既存の leave-one-domain-out の適用である．新しいのは選ぶ対象（参照分類器，重みの付け方，総重み）を K ごとに選ぶ点だけである．
\item 「相補性 J は得の上限」．J（参照が当てて出発点が外す割合 − 逆の割合）は，定義から参照と出発点の正解率の差に等しい．
\end{itemize}

引用の前に原典で確かめる文献が 23 件残っている．

\section{現在の実験の不備}

\subsection{TAC 版の模擬査読の指摘の状態}

TAC 版の模擬査読（機械学習の査読者と感情計算の査読者の 2 人）の指摘 49 件を，関わる筋ごとに数えると 58 件になる．2026-11-01 の照合（\texttt{DAI\_paper/reviews/tac-round1-status.md}）では次の状態だった．

\begin{table}[h]
\centering\small
\caption{模擬査読の指摘の状態（筋ごとに数えた 58 件）}
\begin{tabularx}{\linewidth}{@{}Lrrrrrr@{}}
\toprule
筋 & (a) 対応済み & (b) 比較の設計で対応 & (c) 書き方で対応 & (d) 新しい計算が要る & (e) 直せない & 計 \\
\midrule
1 利得の地図 & 1 & 0 & 0 & 5 & 0 & 6 \\
2 o と d & 3 & 0 & 1 & 4 & 0 & 8 \\
3 ソースの価値 & 0 & 0 & 1 & 5 & 0 & 6 \\
4 ターゲットの価値 & 0 & 2 & 0 & 3 & 0 & 5 \\
5 蒸留と比較 & 0 & 7 & 3 & 5 & 2 & 17 \\
統計・データ・再現・倫理 & 0 & 0 & 7 & 5 & 0 & 12 \\
TAC 版の書き方だけ & 0 & 0 & 4 & 0 & 0 & 4 \\
\midrule
計 & 4 & 9 & 16 & 27 & 2 & 58 \\
\bottomrule
\end{tabularx}
\end{table}

(d) の 27 件は 3 つの実験に振り分けた．保存済みの結果から出せるものは E75（CPU．計算は終わった），学び直しが要るものは E76（GPU．計算中），比較に関わるものは E77〜E79 である．

\subsection{E75 で分かった不備（直した）}

E75 は 13 項目を 394 秒で計算し，検算 550 件に外れは無かった．分かった不備と直した言い方は表 1 の右の列のとおりである．まとめると次の 4 つになる．

\begin{enumerate}
\item 地図の 95\% 区間は 30 回の反復を独立とみなしていた．反復の間でテストの話者が重なるので，区間が狭すぎた．
\item 「得の 5〜6 割は o の直し」という割合は，誤差の大きい模型から出していた．
\item 式 (2) の当たり具合を，模型の中の量でだけ確かめていた．
\item q を「有効次元」と呼んでいたが，ターゲットごとの値が 4.6〜90.7 と散らばり，次元の意味を持たない．
\end{enumerate}

\subsection{計算中の実験で直す不備}

\begin{itemize}
\item 地図はテストの発話ごとの予測を保存していなかったので，話者を単位にした区間と性別ごとの値を出せない．E76 が予測を保存して学び直している．
\item 筋 2 の中心合わせの手の，発話ごとの点数が無い．E76 で保存する．
\item 比較の手法の向きの損を測っていない．E77・E78 が全手法の確率と分類器の係数（重み W と切片 b）を保存し，E79 で測る．
\end{itemize}

\subsection{計算しても直らない不備}

\begin{itemize}
\item 評価に使う 7 群は，すべて手法の設計の段で見ている．このため，要旨と結論の「上回る」は「設計に使ったターゲットで上回った」と限って書く．見ていないコーパス（MSP-IMPROV など）での確認は無い．取得には申請が要る．
\item 第 2 のエンコーダ（HuBERT）での再現は，同じターゲットと同じ話者で測ったので，見ていないデータでの確認にはならない．
\item oursx は ours の結果を見てから足した版である．
\end{itemize}

\subsection{計算中の比較実験で残る弱点}

計算量のため，夜の間に委任で次のように決めた．朝の報告の 4 節で承認を求めている．

\begin{itemize}
\item DANN と MME は，設定を種 0 だけで選び，選んだ設定だけを種 3 つで回す．融合の生徒では DANN と MME を回さない．
\item DANN の最適化（Adam）はミニバッチの組み方で結果が変わるので，種 3 つの平均で読む．
\item Co-teaching は，小さいターゲット（AFEW など）では最適化が終わりきらない．候補の重みを 0 にしても基準の解に届かず，予測の食い違いが 4.4\% 残った．小さいターゲットでは，Co-teaching と基準の差に最適化の不足が混ざる．
\item 合併とターゲットのみの参照の行は，K=5 と 25 だけで回す．
\end{itemize}

\section{比較の不足}

\subsection{今の比較（E61）に無かったもの}

今の比較（E61．2026-10-24 確定）は，基準，ours，oursx，温度つき蒸留，一般化蒸留，硬い疑似ラベル，自己学習，エントロピー最小化，共訓練，共正則化，Saerens の事前補正，massbal（予測の合計をそろえる）の 12 手法だった．模擬査読はこれに次の不足を指摘した．

\begin{enumerate}
\item ドメイン適応の標準の手法（DANN，MME，CORAL）が無い．
\item 温度つき蒸留の参照分類器が ours と違うので，得が参照の違いから来るのか重みの付け方から来るのかを分けられない．
\item 部品を 1 つずつ外した実験（除去実験）が無い．
\item 手法ごとに選んだ設定の候補の数と，選択で取りこぼした量を報告していない．
\item o と d のどちらが変わったかを，手法ごとに測っていない．
\end{enumerate}

\subsection{今回足すもの}

比較の設計の第 8 版（\texttt{DAI\_paper/reviews/comparison-design/proposal\_v8.md}）で，表 4 の手法を足した．

\begin{table}[h]
\centering\small
\caption{今回足す比較相手（群は使う情報で分けた）}
\begin{tabularx}{\linewidth}{@{}p{0.27\linewidth}Lp{0.1\linewidth}@{}}
\toprule
群（使う情報） & 足す手法 & 実験 \\
\midrule
A：候補の出力の割合 & （massbal は E61 にある） & — \\
B1：自分の予測と候補の構造 & 情報最大化（SHOT の IM，TIM），T3A，ラベル伝播 & E77 \\
B2：ソースとターゲットの分布の違い & CORAL ＋基準，DANN & E78 \\
B1＋B2 & MME & E78 \\
C：別のモダリティの分類器の答え & ours と同じ参照分類器を使う温度つき蒸留，Hinton の蒸留，xMUDA 型の相互の模倣，Co-teaching（参照の硬いラベルを雑音のあるラベルとして扱う） & E77・E78 \\
参照の行 & ソースのみ，合併，ターゲットのみ，全ラベルで学んだ上限 & E77 \\
除去実験 & oursx から参照分類器・重みの付け方・総重みを 1 つずつ外した版 & E77 \\
対照 & 候補の重みを 0 にした Co-teaching，写像だけの DANN と CORAL & E78 \\
\bottomrule
\end{tabularx}
\end{table}

全手法で，テストの発話ごとの確率と分類器の係数を保存する．後から別の判定を求められても，学び直さずに出せるようにするためである．

\subsection{回さない手法と，その理由}

\begin{tabularx}{\linewidth}{@{}p{0.28\linewidth}L@{}}
\toprule
手法 & 回さない理由（本文に書く） \\
\midrule
話者ごとの正規化 & テストの話者のラベル無しの発話を使う条件になり，候補からテストの話者を除く本研究の設定と使う情報が違う \\
データ拡張に頼る半教師ありの DA（FixMatch 型，AdaMatch，APE，CDAC） & 原論文の変換は画像の変換で，凍結した発話の特徴では同じ変換を作れない．AdaMatch の分布の合わせ込みとしきい値は，ours の釣り合いと硬いラベルのしきい値として比較に含まれる \\
MM-SADA & モダリティの対応を当てる課題は，エンコーダを学ばない設定では使えない．合わせ込みの部分を移すと DANN と同じ形になる \\
C2KD & 教師も更新するので，ソースで固定した参照分類器の設定に合わない \\
部分空間の合わせ込み，線形の最適輸送 & 固定した特徴の上の浅い DA の比較（SKADA-Bench）で単純な線形の写像が頑健だったことを引き，代表として CORAL を回す \\
Amanda & 本文を確かめてから理由を書く（未確認） \\
エンコーダの微調整 & 主張を凍結した特徴に限る．LoRA の実験（E63）を付録に置く \\
\bottomrule
\end{tabularx}

\subsection{範囲の外に置いたもの}

委任で次を範囲の外にした．本文の限界の節に書く．

\begin{itemize}
\item 比較実験を第 2 のエンコーダでやり直すこと．
\item VT と AT の世界の生徒（融合の生徒も単モダリティの生徒も回さない．比較は，全モダリティがそろう AVT と，ターゲットが最も多い AV の世界で行う．VT・AT の融合の生徒には E61 の古い手順の値があり，探索の値として扱う）．
\item 見ていないコーパスでの確認（ステップ 5．データの取得の申請が要る）．
\end{itemize}

\subsection{判定の決め方}

判定の規則は計算の前に E79 の設計に書いた．「上回った」は，1 単位（生徒 × K × コーパス）で差の点の値が事前に決めた幅 δ 以上で，区間が 0 を外れることとする．要旨の文は「21 の生徒 × K のうち n で上回った」の数え上げの形で書き，都合のよい生徒と K だけを取り出さない．偶然に「上回った」になる単位の見込みの数（単位の数 × 0.025）も並べて書く．この設計で見つけられるのは，真の差がおよそ 1.5δ 以上の場合である．

\section{必要な時間}

\subsection{今の見込み}

表 5 に，残りの計算の見込みを示す．時刻は 2026-11-01 15 時の状態である．GPU は 1 枚で，E76 と E77 が分け合っている．2 本より多く並べると GPU のメモリの合計が上限（24 GB）に近づき，計算が 5〜25 倍遅くなる（今朝に起きた）．このため E78 は E77 が終わってから 1 本で回す．

\begin{table}[h]
\centering\small
\caption{残りの計算の見込み（2026-11-01 15 時）}
\begin{tabularx}{\linewidth}{@{}p{0.08\linewidth}p{0.2\linewidth}p{0.2\linewidth}p{0.17\linewidth}L@{}}
\toprule
実験 & 内容 & 状態 & 残りの見込み & 見積もりの根拠 \\
\midrule
E76 & 筋 1〜3 の学び直し（GPU） & 42 ターゲットのうち 23 が終わった & 7〜9 時間（今夜に終わる） & 終わった 23 ターゲットの所要（1 つ 5〜63 分．\texttt{logs/e76.log}） \\
E77 & 比較の第 1 波（GPU） & 850 単位のうち 20 単位が終わった（AVT の世界の IEMOCAP 即興．6.8 時間） & 約 80 時間（3.3 日） & 終わった 20 単位から作った 1 単位の所要の式（下の 6.2 節） \\
E78 & 比較の第 2 波（GPU） & E77 の後に自動で始まる & 3〜11 日（不確か） & 検査役の試走から出した 173〜272 時間を，単独で回すと半分ほどになると見た値．試走は GPU が遅くなっていた時間帯（07〜08 時）に測ったので，大きめに出ている見込みである \\
E79 & 判定の解析（CPU） & E77 と E78 の後 & 2〜4 時間 & 検査役が人工の入力で測った値からの外挿 \\
\bottomrule
\end{tabularx}
\end{table}

今の設定のままだと，E77 が終わるのは 11 月 4〜5 日ごろ，E78 が終わるのは 11 月 7〜16 日ごろである．E79 の判定はその数時間後に出る．E78 の見積もりの幅が広いのは，E78 の 1 単位の所要を正常な GPU の状態で測れていないためである．E77 が終わって E78 が始まると，最初のターゲットの所要から見積もりを測り直せる．

\subsection{E77 の時間の内訳}

E77 の 1 単位は，1 つのターゲット（評価のターゲットか仮のターゲット）の 1 回の反復で，全手法の全設定の候補を解く計算である．終わった 20 単位の実測から，1 単位の所要は「11.4 分＋候補の発話 1,000 本ごとに 1.22 分」で表せた（AVT の世界．候補は IEMOCAP 即興で約 2,350 本，MOSEI で約 15,100 本，MELD で約 9,200 本）．この式は IEMOCAP 即興の 20 単位の合計 6.8 時間を再現する（式の値 6.75 時間）．

時間の大部分は仮のターゲットの計算である．IEMOCAP 即興では，評価の 10 反復に 2.2 時間，仮のターゲット（MOSEI と MELD を各 5 回）に 3.8 時間かかった．仮のターゲットでは，すべての手法のすべての設定の候補を解いて，その中から設定を選ぶからである．

手法ごとに見ると，共正則化が実時間 6.8 時間のうち 2.4 時間（35\%）を使っていた（表 6）．共正則化は 1 単位あたり約 7 分かかり，この時間は候補の数にほとんど依らない．

\begin{table}[h]
\centering\small
\caption{E77 の AVT の世界の IEMOCAP 即興での手法ごとの所要（\texttt{results/e77\_fuse-audio-video-text/IEMOCAP-impro\_conv.csv}．割合は実時間 6.8 時間（24,524 秒．\texttt{logs/e77b\_fuse-audio-video-text.log}）に対する値）}
\begin{tabular}{@{}lrrrr@{}}
\toprule
手法 & 評価（時間） & 仮のターゲット（時間） & 合計（時間） & 割合 \\
\midrule
共正則化 & 1.18 & 1.22 & 2.40 & 35\% \\
ours・oursx と除去実験の弱いラベルの行 & 0.15 & 0.65 & 0.80 & 12\% \\
ラベル伝播 & 0.10 & 0.43 & 0.53 & 8\% \\
エントロピー最小化 & 0.08 & 0.43 & 0.51 & 7\% \\
自己学習 & 0.11 & 0.36 & 0.47 & 7\% \\
共訓練 & 0.13 & 0.23 & 0.36 & 5\% \\
参照の行 & 0.29 & — & 0.29 & 4\% \\
温度つき蒸留 & 0.04 & 0.20 & 0.24 & 4\% \\
\bottomrule
\end{tabular}
\end{table}

\note{残りの約 1.2 時間は，記録の区分を持たない計算（情報最大化，T3A，共訓練の内側の解き直し）と，データの読み込みと保存である．}



\section{削減案}

\begin{tcolorbox}
\textbf{案}：E77 の残りと E78 で，結果の値を変えない省き方を 1 つ（下の 1），比較の判定に使わない行の省き方を 1 つ（下の 2）入れる．E77 の残りは約 80 時間から約 58 時間に減り，E77 は約 1 日早く終わる見込みである．E78 も 1 の分（1 割ほど）短くなる．
\end{tcolorbox}

\subsection{1．IEMOCAP 台本の仮のターゲットの計算を，即興の結果の写しにする}

\begin{itemize}
\item \textbf{何をするか}：IEMOCAP 台本を評価のターゲットにしたときの仮のターゲットの計算をやめ，IEMOCAP 即興のときに計算した同じ単位の結果を写して使う．
\item \textbf{理由}：仮のターゲットの計算に使うものは，即興と台本で全部同じである．即興と台本は同じ話者なので，どちらがターゲットでも，もう一方はソースから外れる．そのため使うソースの集合（AVT の世界なら MOSEI と MELD）が同じになる．出発点の 5 つの設定値も，5 種類の生徒すべてで一致することを確かめた．データの分け方も同じである．つまり今は同じ計算を 2 回している．説明資料 p5 でも，即興と台本で選ばれた設定が完全に同じだった．
\item \textbf{確かめ方}：今走っている E77 の台本のターゲットの最初の 1 単位（MOSEI の反復 0）が，即興の同じ単位の出力と一致することを確かめてから写す．
\item \textbf{減る時間}：E77 の残りで約 11 時間，E78 で 1 割ほど．
\item \textbf{失うもの}：無い．
\end{itemize}

\subsection{2．共正則化は仮のターゲットで解かず，評価のターゲットでは相手に有利な値で比べる}

\begin{itemize}
\item \textbf{何をするか}：共正則化だけ，仮のターゲットでは計算しない．評価のターゲットでは今どおり 4 つの設定を全部解き，テストのラベルで選んだ最良の値（オラクルの値）と oursx を比べる．
\item \textbf{理由}：共正則化は E77 の実時間の 35\% を使う．比較の設計では群 C の探索の行で，筋 4 と筋 5 の判定（P4，P5-1〜P5-4）には使わない．テストで選んだ最良の値は，仮のターゲットで選んだ値より必ず高いか等しい．このため oursx がオラクルの値を上回れば，仮のターゲットで選んだ値も上回ったと言える．相手を不利にする向きの偏りは入らない．
\item \textbf{減る時間}：E77 の残りで約 13 時間（AVT の世界で約 3.5 時間，AV の世界で約 9.4 時間）．
\item \textbf{失うもの}：共正則化の「仮のターゲットで選んだ値」と「選択の取りこぼし」の列．oursx がオラクルの値を上回らなかった単位では，仮のターゲットで選んだ値との比較について何も言えない．E61 では共正則化の UAR は基準とほぼ同じか下だった．
\end{itemize}

\subsection{進め方}

\begin{enumerate}
\item 台本を直す（1 と 2 を引数で選べるようにする）．E79 の設計の群 C の探索の行を「共正則化はオラクルの値」と書き直す．
\item 検査役（gate-code，gate-compare）に差分を検査させる．
\item 今走っている E77 は，台本のターゲットの仮のターゲットを計算している途中である．これを止めて，直した台本で台本のターゲットから流し直す．止めて失うのは，写しで置き換わる計算だけである．E78 の待ち受けはそのまま使える．
\end{enumerate}

1 と 2 の作業と検査に 1〜2 時間かかる．その間も E77 は今の設定のまま進む．


\end{document}
